\chapter*{Kết luận: Bình minh của AI tạo sinh trong học thuật}
\addcontentsline{toc}{chapter}{Kết luận: Bình minh của AI tạo sinh trong học thuật}
\markboth{Kết luận: Bình minh của AI tạo sinh trong học thuật}{Kết luận: Bình minh của AI tạo sinh trong học thuật}

\subsection{Tóm tắt}

Chương kết luận này suy ngẫm về tác động mang tính chuyển đổi của AI tạo sinh đối với nghiên cứu học thuật, viết học thuật, giảng dạy và thực tiễn thể chế. Thay vì xem AI chỉ như một công cụ nâng cao hiệu suất, chương xem xét cách AI đang định hình lại những khía cạnh nền tảng của công việc học thuật, từ phương pháp nghiên cứu đến quyền tác giả, phản biện đồng cấp (peer review) và phương pháp sư phạm. Chương lập luận rằng AI tạo sinh đại diện cho một sự chuyển đổi mang tính hệ hình trong quá trình sản xuất tri thức, đòi hỏi các học giả phải suy nghĩ lại các chuẩn mực truyền thống về lao động trí tuệ, tính nguyên bản và ranh giới giữa các ngành. Chương khám phá cả những cơ hội lẫn những thách thức đạo đức đi kèm với việc tích hợp AI vào môi trường học thuật, bao gồm các vấn đề về thiên kiến, công bằng, tính bền vững và liêm chính học thuật. Nhấn mạnh khái niệm AI cộng tác, chương ủng hộ một tầm nhìn trong đó sự sáng tạo của con người và phán đoán đạo đức cùng vận hành song hành với trí tuệ máy móc nhằm thúc đẩy học thuật. Chương kết thúc bằng lời kêu gọi sự tham gia mang tính phản biện, trách nhiệm của các thể chế và sự suy ngẫm tập thể, thúc giục giới học thuật định hình tương lai của việc tích hợp AI theo cách củng cố, thay vì làm tổn hại, sứ mệnh cốt lõi của mình.

\subsection{Từ khóa}

AI tạo sinh, chuyển đổi học thuật, phương pháp luận nghiên cứu, quyền tác giả học thuật, liêm chính học thuật, đạo đức AI, AI cộng tác, nhận thức luận, nghiên cứu liên ngành, xuất bản học thuật, chính sách giáo dục đại học, năng lực AI (AI literacy), quản trị học thuật, tính bền vững, thay đổi công nghệ

Giới học thuật hiện nay đồng thời đón nhận và đặt dấu hỏi về những công cụ mạnh mẽ mới này. Các nhà nghiên cứu từng mất hàng tuần để thực hiện tổng quan tài liệu nay có thể tạo ra những bản tổng hợp toàn diện chỉ trong vài giờ. Các học giả gặp khó khăn với tình trạng bí ý tưởng khi viết có thể tương tác với các hệ thống AI gợi ý những cách diễn đạt, cấu trúc và lập luận thay thế. Sinh viên đang thử nghiệm những hình thức học tập và kiến tạo tri thức mới nhờ các hệ thống gia sư thông minh và cơ chế phản hồi cá nhân hóa. Tuy nhiên, song hành với những năng lực đáng kinh ngạc ấy là những câu hỏi sâu sắc về bản chất của quyền tác giả, nền tảng của liêm chính học thuật, việc tiếp cận công bằng các lợi thế công nghệ, và tương lai của lao động cũng như chuyên môn học thuật.

Bước ngoặt công nghệ này đòi hỏi không chỉ sự thích ứng thực tiễn mà còn sự suy ngẫm sâu sắc về các giá trị và mục đích thúc đẩy công việc học thuật. Khi AI tạo sinh ngày càng thấm vào mọi khía cạnh của thực hành học thuật, từ khi hình thành câu hỏi nghiên cứu đến các giai đoạn cuối của xuất bản và phổ biến, chúng ta phải cùng nhau xác định những công cụ này có thể phục vụ tốt nhất sứ mệnh nền tảng của giới học thuật như thế nào: thúc đẩy tri thức và sự hiểu biết của nhân loại thông qua hoạt động tìm hiểu nghiêm cẩn, sáng tạo và có đạo đức.

Trọng tâm của quá trình tiến hóa này là sự xuất hiện của AI cộng tác, những hệ thống được thiết kế không phải để thay thế các học giả mà để làm việc song hành với họ như những đối tác trí tuệ. Các khung cộng tác này thể hiện một sự chuyển dịch căn bản: từ việc xem AI chỉ như công cụ tự động sang việc hiểu chúng như những tác nhân tham gia vào quá trình kiến tạo tri thức. Thông qua các giao diện tinh vi hỗ trợ tương tác theo thời gian thực, tinh chỉnh ý tưởng lặp đi lặp lại và các thế mạnh nhận thức bổ sung cho nhau, các hệ thống AI cộng tác đang định hình lại chính bản chất của làm việc nhóm trong học thuật. Chúng mở ra những hình thức hợp tác trí tuệ mới, nơi sự sáng tạo, hiểu biết bối cảnh và phán đoán đạo đức của con người kết hợp với sức mạnh tính toán, khả năng nhận dạng mẫu và xử lý thông tin ở quy mô lớn.

Trong những trang tiếp theo, chúng tôi tổng hợp các nhận định chính từ quá trình khảo sát AI tạo sinh trong bối cảnh học thuật, đưa ra những kết luận đề cập cả các hệ quả thực tiễn trước mắt lẫn những câu hỏi triết học sâu xa hơn mà cuộc cách mạng công nghệ này đặt ra. Mục đích của chúng tôi không phải là đưa ra câu trả lời dứt khoát, quả thực, nhiều câu hỏi quan trọng vẫn còn bỏ ngỏ khi các công nghệ này tiếp tục phát triển nhanh chóng, mà là xây dựng một nền tảng thấu đáo cho đối thoại và cân nhắc liên tục trong cộng đồng học thuật. Nhờ trí tuệ tập thể như vậy, chúng ta có thể vượt qua thời khắc chuyển mình này theo cách khai thác tiềm năng phi thường của AI đồng thời gìn giữ những yếu tố mang đậm tính người vốn mang lại cho công việc học thuật ý nghĩa và giá trị tối hậu.

\section{Sự chuyển đổi mô hình trong phương pháp luận nghiên cứu}

AI tạo sinh không chỉ là một bước tiến công nghệ mang tính gia tăng, mà là một sự chuyển đổi mô hình căn bản trong cách nghiên cứu học thuật được hình thành về mặt khái niệm và được tiến hành. Từ những giai đoạn đầu của việc hình thành ý tưởng cho đến khâu phổ biến kết quả cuối cùng, các công cụ AI đang định nghĩa lại quy trình nghiên cứu truyền thống, đẩy nhanh các quá trình khám phá và nâng cao đáng kể năng lực trí tuệ của con người. Ngày nay, công nghệ này cho phép các nhà nghiên cứu xử lý và tổng hợp những mảng tài liệu rộng lớn mà trước đây phải mất hàng tháng hoặc hàng năm nỗ lực của con người, đề xuất các giả thuyết mới bằng cách nhận diện các mẫu hình xuyên qua những lĩnh vực dường như không liên quan, và quản lý các tập dữ liệu phức tạp, đa chiều với hiệu quả và chiều sâu hiểu biết chưa từng có.

Sự chuyển dịch này đặc biệt rõ nét trong cách các học giả hiện nay tiếp cận tổng quan tài liệu và tổng hợp tri thức. Trong khi các thế hệ nhà nghiên cứu trước bị giới hạn bởi những hạn chế về nhận thức khi xử lý khối lượng tài liệu học thuật tăng theo cấp số nhân, thì các công cụ AI tạo sinh nay có thể nhanh chóng quét, phân loại và trích xuất những hiểu biết then chốt từ hàng nghìn bài báo, giúp hiểu các lĩnh vực nghiên cứu một cách toàn diện và tinh tế hơn. Tương tự, trong phân tích dữ liệu, khả năng của AI trong việc nhận diện các mẫu hình và mối tương quan tinh vi đang mở rộng ranh giới của những câu hỏi có thể được đặt ra và giải đáp một cách có ý nghĩa thông qua điều tra thực nghiệm.

Những hệ quả này vượt ra ngoài việc chỉ gia tăng hiệu quả. Khi AI ngày càng được tích hợp vào quá trình nghiên cứu, chúng ta có thể chứng kiến những thay đổi căn bản trong cách tri thức được kiến tạo và thẩm định xuyên suốt các ngành, có khả năng thách thức các khung nhận thức luận đã tồn tại lâu đời và các lối chính thống về phương pháp luận. Điều này gợi ý không chỉ những thay đổi mang tính thực tiễn trong thực hành nghiên cứu mà còn những câu hỏi triết học sâu sắc hơn về bản chất của việc sản xuất tri thức trong thời đại cộng tác giữa con người và máy móc.

\section{Chuyển đổi trong viết và xuất bản học thuật}

Tác động của AI đối với việc viết và xuất bản học thuật cũng là một sự chuyển đổi sâu sắc đối với các thực hành cốt lõi của giới học thuật. Trong khi các trợ lý AI tạo sinh (generative AI) có thể giúp cấu trúc, soạn thảo và trau chuốt văn bản học thuật, chúng cũng đặt ra những câu hỏi mang tính then chốt về việc ghi nhận quyền tác giả, đóng góp trí tuệ và chính bản chất của giao tiếp học thuật. Sự phát triển của các công cụ này gợi ý về một mối quan hệ cộng sinh đang hình thành giữa sự sáng tạo của con người và năng lực của máy móc, mối quan hệ mà rốt cuộc có thể định nghĩa lại ý nghĩa của việc trở thành một tác giả trong các bối cảnh học thuật.

Khả năng của các hệ thống AI trong việc tạo ra văn phong học thuật mạch lạc, có cấu trúc tốt đặt ra những câu hỏi quan trọng về các kỹ năng sẽ được coi trọng ở các học giả tương lai. Nếu các khía cạnh cơ học của việc viết ngày càng được tự động hóa, thì sự nổi bật trong học thuật có thể ngày càng phụ thuộc vào đổi mới khái niệm, tầm nhìn lý thuyết và sự tổng hợp sáng tạo, thay vì năng lực viết truyền thống. Sự thay đổi này có thể đặc biệt ảnh hưởng đến những truyền thống học thuật mà ở đó các phong cách viết và quy ước tu từ riêng biệt từ lâu đã đóng vai trò là dấu hiệu của bản sắc và sự thuộc về học thuật.

Ngoài ra, các công cụ AI đang định hình lại chính bối cảnh xuất bản, từ quản lý trích dẫn tự động đến các hệ thống phát hiện đạo văn tinh vi và phản biện đồng cấp (peer review) có hỗ trợ của AI. Những diễn biến này gợi ý một tương lai trong đó toàn bộ hệ sinh thái xuất bản ngày càng được trung gian hóa bởi các công nghệ AI, có khả năng giải quyết những bất cập tồn tại lâu nay, đồng thời đưa vào những phức tạp và thách thức mới đối với các chức năng gác cổng học thuật truyền thống.

\section{Những yêu cầu đạo đức và thách thức quản trị}

Các khía cạnh đạo đức của việc ứng dụng AI trong học thuật không thể bị xem nhẹ và đòi hỏi sự quan tâm khẩn thiết, bền bỉ từ cộng đồng học giả. Khi các công nghệ này lan rộng trong các cơ sở giáo dục, các học giả cá nhân, nhà quản lý và nhà hoạch định chính sách phải đối mặt với những câu hỏi ngày càng phức tạp liên quan đến bảo vệ dữ liệu, quyền riêng tư, quyền sở hữu trí tuệ, thiên lệch thuật toán và liêm chính học thuật. Việc xây dựng các khung đạo đức và cơ chế quản trị vững chắc sẽ là điều thiết yếu để tích hợp AI một cách có trách nhiệm, phù hợp với các giá trị cốt lõi của học thuật.

Đặc biệt đáng lo ngại là các vấn đề về công bằng và khả năng tiếp cận. Nếu các công cụ AI tiên tiến chủ yếu chỉ dành cho những cơ sở và học giả có nhiều nguồn lực, chúng có nguy cơ làm trầm trọng thêm những bất bình đẳng sẵn có trong sản xuất tri thức và ảnh hưởng học thuật. Tương tự, tác động môi trường của việc huấn luyện và triển khai các mô hình ngôn ngữ lớn cùng các hệ thống AI khác đặt ra những câu hỏi về tính bền vững và trách nhiệm trong hạ tầng nghiên cứu.

Các mối quan ngại về liêm chính học thuật cũng cấp bách không kém. Khi văn bản do AI tạo ra ngày càng tinh vi và khó phân biệt với văn bản của con người, những quan niệm truyền thống về đạo văn và việc đánh giá kết quả học tập của sinh viên đòi hỏi phải được xem xét lại một cách căn bản. Các cơ sở giáo dục phải xây dựng những cách tiếp cận mới để đánh giá đóng góp học thuật và thành tích của sinh viên, thừa nhận thực tế về sự hỗ trợ của AI đồng thời vẫn duy trì các chuẩn mực có ý nghĩa về sự trung thực và tính nghiêm cẩn trí tuệ.

Hơn nữa, việc dựa vào các hệ thống AI được huấn luyện chủ yếu trên văn liệu học thuật phương Tây có nguy cơ củng cố những thiên kiến sẵn có trong sản xuất tri thức và làm lu mờ các truyền thống trí tuệ cũng như nhận thức luận phi phương Tây. Giải quyết những mối quan ngại này không chỉ cần các giải pháp kỹ thuật mà còn cần sự dấn thân sâu hơn vào các câu hỏi về việc tri thức của ai được coi trọng và làm thế nào để các truyền thống trí tuệ đa dạng được thể hiện một cách có ý nghĩa trong học thuật có AI hỗ trợ.

\section{Các quỹ đạo tương lai và sự tiến hóa của học thuật}

Quỹ đạo phát triển của AI cho thấy rất rõ rằng ảnh hưởng của nó trong giới học thuật sẽ tiếp tục sâu sắc và mở rộng hơn trong những năm tới. Như đã nêu bật trong phần khám phá các xu hướng tương lai của cuốn sách, chúng ta đang đứng ở thời điểm khởi đầu của một giai đoạn chuyển đổi, nhiều khả năng sẽ định nghĩa lại thực hành học thuật trên mọi lĩnh vực theo những cách vừa có thể lường trước vừa không thể dự đoán. Những học giả phát triển được sự thành thạo với các công nghệ này có thể đạt được lợi thế đáng kể về năng suất nghiên cứu, tác động và ảnh hưởng, trong khi những người từ chối áp dụng có nguy cơ bị gạt ra bên lề trong các môi trường học thuật ngày càng cạnh tranh.

Việc tích hợp các hệ thống AI đa phương thức có khả năng làm việc xuyên suốt văn bản, hình ảnh, mã nguồn và các loại dữ liệu khác hứa hẹn mang lại những thay đổi đặc biệt sâu sắc đối với nghiên cứu liên ngành. Những công cụ này có thể mở ra các hình thức tổng hợp tri thức mới vượt qua ranh giới truyền thống giữa các ngành và những chia cắt về phương pháp luận, từ đó có khả năng thúc đẩy các đột phá trong những lĩnh vực vấn đề phức tạp, từ biến đổi khí hậu đến sức khỏe cộng đồng.

Các ứng dụng mới nổi của AI tạo sinh trong bối cảnh giáo dục cũng gợi ý về những thay đổi sâu rộng đối với thực hành giảng dạy và trải nghiệm học tập của sinh viên. Từ các lộ trình học tập cá nhân hóa đến các môi trường mô phỏng tinh vi, những công nghệ này có thể biến đổi cách tri thức chuyên ngành và kỹ năng nghiên cứu được truyền đạt cho các thế hệ học giả mới, với những hàm ý quan trọng đối với thiết kế chương trình đào tạo, các phương pháp sư phạm và đánh giá giáo dục.

Hơn nữa, bản thân tốc độ phát triển ngày càng nhanh của AI đặt ra thách thức cho các cấu trúc quản trị và chính sách học thuật, vốn thường tiến hóa chậm hơn nhiều so với các công nghệ mà chúng tìm cách điều chỉnh. Sự lệch pha về thời gian này cho thấy cần có những cách tiếp cận dựa trên nguyên tắc và thích ứng hơn đối với quản trị công nghệ trong học thuật, có thể đáp ứng những năng lực thay đổi nhanh chóng đồng thời duy trì các giá trị cốt lõi của thể chế.

\section{Tham gia có suy xét và trách nhiệm của cộng đồng}

Cuộc cách mạng công nghệ này đòi hỏi sự tham gia có suy xét, thay vì chấp nhận thiếu phê phán hay bác bỏ hoàn toàn. Cộng đồng học thuật phải tham gia một cách chủ động và phê phán vào việc định hình cách các công cụ AI phát triển và hòa nhập vào thực hành học thuật, bảo đảm chúng củng cố thay vì làm suy yếu những giá trị cốt lõi của hoạt động nghiên cứu học thuật: tính chặt chẽ, tính minh bạch, liêm chính và sự thăng tiến của tri thức nhân loại.

Đặc biệt quan trọng là giữ gìn không gian cho những dạng thức sáng tạo, thấu hiểu và phán đoán riêng biệt của con người vốn vẫn nằm ngoài khả năng của AI. Trong khi đón nhận những hiệu quả và khả năng mới mà các công cụ AI mang lại, học giả phải duy trì nhận thức phê phán về những hạn chế và thiên kiến của chúng, và sử dụng chúng như phần bổ trợ cho, chứ không phải thay thế, lao động trí tuệ của con người.

Các phản ứng mang tính thể chế sẽ có vai trò then chốt trong việc vượt qua giai đoạn chuyển đổi này. Các trường đại học, nhà tài trợ nghiên cứu và hội học thuật phải xây dựng những chính sách chu đáo, khuyến khích đổi mới có trách nhiệm đồng thời bảo vệ các giá trị học thuật và giải quyết những mối quan ngại về đạo đức. Điều này bao gồm đầu tư vào năng lực hiểu biết về AI cho giảng viên và sinh viên, xây dựng các cấu trúc quản trị phù hợp, và chủ động tham gia vào các cuộc thảo luận xã hội rộng lớn hơn về quy định và phát triển AI.

Suy cho cùng, việc AI tạo sinh định hình lại giới học thuật như thế nào sẽ không phụ thuộc vào thuyết tất định công nghệ mà phụ thuộc vào những lựa chọn tập thể của cộng đồng học thuật. Bằng cách tiếp cận các công nghệ này với cả sự cởi mở trước tiềm năng chuyển đổi của chúng lẫn nhận thức phê phán về những hạn chế và rủi ro, giới học thuật có thể khai thác AI như một công cụ mạnh mẽ để thúc đẩy tri thức nhân loại, đồng thời giữ gìn những yếu tố nhân văn thiết yếu của hoạt động tìm hiểu trí tuệ.
